% Subsection 5.2.1: Base 与 LoRA 实验结果

本节汇总 4 组训练侧对照（H1、H4、H7、H9，推理侧固定为 Direct）的 pass@$k$。所有数值按 4.1 节表~\ref{tab:ch4-matrix} 结果命名规则从 \texttt{evaluation/results/} 读取。

\subsubsection{主结果表}
表~\ref{tab:ch5-direct} 给出 Direct 推理侧下四种训练配置的 pass@$k$，三个 $k$ 档同时给出以便观察「单次命中」与「预算内命中」之间的差距。

\begin{table}[!htb]
  \centering
  \caption{Direct 推理侧下四种训练配置的 pass@$k$（单位：\%）}
  \label{tab:ch5-direct}
  \footnotesize
  \begin{tabular}{lllccc}
    \toprule
    编号 & 训练侧 & 推理侧 & pass@1 & pass@5 & pass@10 \\
    \midrule
    H1 & Base               & Direct & 52.40 & 67.10 & 72.80 \\
    H4 & LoRA\_code         & Direct & 60.50 & 73.80 & 78.90 \\
    H7 & LoRA\_cot\_zs      & Direct & 57.20 & 70.80 & 76.40 \\
    H9 & LoRA\_cot\_fs($k$) & Direct & 58.50 & 71.90 & 77.30 \\
    \bottomrule
  \end{tabular}
\end{table}

\subsubsection{观察要点}
表~\ref{tab:ch5-direct} 在主干上对应以下三条差分，具体数值分析延至 5.3.1 节：
\begin{itemize}
    \item \textbf{H4 − H1}：纯代码 LoRA 在 Direct 推理下的训练侧主效应；
    \item \textbf{H7 − H4}：将 CoT 格式注入训练监督后，相对仅代码监督的增益；
    \item \textbf{H9 − H7}：在 CoT 监督基础上进一步注入 $k$ 条训练 prompt 示例后的独立增益。
\end{itemize}
由表~\ref{tab:ch5-direct} 可见，训练侧从 Base 依次过渡到 \texttt{LoRA\_code}、\texttt{LoRA\_cot\_zs}、\texttt{LoRA\_cot\_fs($k$)} 时，pass@1 走势为\textbf{先升（+8.1\,pp）后降（−3.3\,pp）再回升（+1.3\,pp）}：\texttt{LoRA\_code} 在 Direct 推理下取得最大值 60.50\%，\texttt{LoRA\_cot\_zs} 在同一推理侧反而回落到 57.20\%，加入训练侧 few-shot 后 \texttt{LoRA\_cot\_fs($k$)} 略有恢复但仍未追平 \texttt{LoRA\_code}；pass@5/10 走势完全同号。该「单峰」走势是 5.3.1 节讨论 LoRA 有效性以及 CoT 监督在 Direct 推理下出现「格式开销」的起点。

\subsubsection{数据自检}
除主结果外，每份 Direct 评测 JSON 都记录了原始 \texttt{solution} 在本运行环境下的通过率，四组间应保持一致且接近 100\%；若出现显著偏离，则说明该次运行的执行器或依赖发生变化，需要回滚环境后重跑，避免把环境扰动误计为方法差异。